\documentclass[12pt]{article}
\usepackage[utf8]{inputenc}
\usepackage[T1]{fontenc}
\usepackage{hyperref}
\title{Korlai Portuguese Creole Language Resources for K-12 Curriculum Development}
\author{Compiled by Claude Sonnet 5 (claude-sonnet-5) for Andr\'as Kornai}
\date{}
\begin{document}
\maketitle
\section{Introduction}
Korlai Portuguese Creole (ISO 639-3 vkp), locally called Christi or Kristi, is a Portuguese-based creole spoken in Korlai village near Chaul, coastal Maharashtra, India. Estimates are under 1,000 speakers, with one source giving about 800. It is under pressure from Marathi. The main scholarly source is J. Clancy Clements' 1996 monograph.
\section{Primary Educational Resources}
None found. No school materials in Korlai Creole were seen.
\section{Linguistic and Reference Materials}
\begin{itemize}
\item Library record for Clements, \emph{The Genesis of a Language: The Formation and Development of Korlai Portuguese} (Benjamins, 1996): \url{https://libcat.mcphs.edu/Record/in00000317775}
\item APiCS survey of Korlai: \url{https://apics-online.info/surveys/40}
\item Sahapedia article on Korlai Creole: \url{https://www.sahapedia.org/korlai-creole-history-and-characteristics-indo-portuguese-language-spoken-coastal-maharashtra}
\item LINGUIST List 7.818: \url{https://linguistlist.org/issues/7/818/}
\end{itemize}
\section{Religious and Community Resources}
No community organisation or religious materials were found.
\section{Digital and Audio Resources}
Popular press coverage: \url{https://www.thebridgechronicle.com/amp/story/lifestyle/art-culture/have-you-heard-of-the-portuguese-speaking-community-in-maharashtra}. No audio corpus was found in the results.
\section{Curriculum Development Context}
The language has no established school use. A curriculum would need a community-agreed orthography, with Marathi as the dominant surrounding language, and could build on Clements' descriptive work.
\end{document}